\answer{ex:18:1}{(a)~$0.9949\le w\le1.0049$, soit $|1-w|\lesssim0.005$. (b)~$K\ge400$ synapses.}
\answer{ex:18:2}{$(\omega_R\tau_w)^2=\sqrt{\alpha(\alpha+2+2\varrho)}/\varrho-1$, $\varrho=\tau_m/\tau_w$; $f_R\approx11.1$~Hz; $|Z(\omega_R)|/|Z(0)|=4.6$.}
\answer{ex:18:3}{(a)~$f=1/\bigl(\tau\ln\frac{\mu}{\mu-\theta}\bigr)$; pour $1$~Hz, il faut $\mu/\theta-1\approx3.7\cdot10^{-44}$. (b)~$T=\pi\tau/\sqrt\mu$, $f\approx3.2$~Hz: $f\propto\sqrt\mu$.}
\answer{ex:18:4}{L'intégrateur décharge pour $f\lesssim17.7$~Hz (filtre passe-bas), le résonateur seulement dans la bande $5.5$--$22$~Hz (filtre passe-bande).}
\answer{ex:18:5}{(a)~$T_{\min}=\frac{\tau}{w-1}\ln\frac{0.5}{0.3}\approx10.2\ms$. (b)~$B>w-1-0.2=0.3$.}
\answer{ex:18:6}{Constante de temps du réglage $\tau/(\eta\langle r^2\rangle)$, $\eta=\tau\ln10/60\,\text{s}\approx3.8\cdot10^{-3}$; pour $I\neq0$, le poids se décale: il faut retrancher de $e$ une copie de la commande $I/\tau$.}
\answer{ex:18:7}{Question ouverte. Le résonateur ne gagne qu'un facteur $1.35$ à $11$~Hz et perd un facteur $17$ à $1$~Hz; le principal avantage de la reconfiguration est le choix de la bande par le seuil (exercice~\ref{ex:18:4}).}
